\documentclass[11pt, a4paper]{article}

% bfw-style.tex — gemeinsame Style-Definitionen für alle Handouts/Aufgabenhefte
% Voraussetzung: Kompilation mit xelatex (wegen fontspec)

\usepackage[ngerman]{babel}
\usepackage{geometry}
\geometry{a4paper, top=2.2cm, bottom=2.2cm, left=2.3cm, right=2.3cm}

\usepackage{fontspec}
% Fallback-Kette: Source Sans Pro -> Calibri -> Segoe UI -> Latin Modern Sans
% (Latin Modern ist immer mit MiKTeX vorhanden)
\IfFontExistsTF{Source Sans Pro}{%
  \setmainfont{Source Sans Pro}[Ligatures=TeX]%
}{\IfFontExistsTF{Calibri}{%
  \setmainfont{Calibri}[Ligatures=TeX]%
}{\IfFontExistsTF{Segoe UI}{%
  \setmainfont{Segoe UI}[Ligatures=TeX]%
}{%
  \setmainfont{Latin Modern Sans}[Ligatures=TeX]%
}}}
\IfFontExistsTF{Source Code Pro}{%
  \setmonofont{Source Code Pro}[Scale=0.92]%
}{\IfFontExistsTF{Consolas}{%
  \setmonofont{Consolas}[Scale=0.92]%
}{%
  \setmonofont{Latin Modern Mono}[Scale=0.92]%
}}

\usepackage{xcolor}
% Farbpalette: hell, freundlich, modern
\definecolor{bfwPrimary}{HTML}{0EA5A4}    % Petrol/Türkis - Primär
\definecolor{bfwPrimaryDark}{HTML}{0F766E} % dunkler Petrol für Headings
\definecolor{bfwAccent}{HTML}{F59E0B}     % Amber - Akzent
\definecolor{bfwSuccess}{HTML}{10B981}    % Grün - Erfolg / Lösung
\definecolor{bfwWarn}{HTML}{F97316}       % Orange - Achtung
\definecolor{bfwInk}{HTML}{0F172A}        % fast Schwarz
\definecolor{bfwInkSoft}{HTML}{475569}    % gedämpftes Grau
\definecolor{bfwBgSoft}{HTML}{F8FAFC}     % off-white
\definecolor{bfwBgTeal}{HTML}{ECFEFF}     % helles Türkis
\definecolor{bfwBgAmber}{HTML}{FEF3C7}    % helles Gelb
\definecolor{bfwBgRose}{HTML}{FFE4E6}     % helles Rosé für Eselsbrücken

\usepackage[colorlinks=true, linkcolor=bfwPrimaryDark, urlcolor=bfwPrimaryDark]{hyperref}
\usepackage{enumitem}
\setlist[itemize]{leftmargin=*, itemsep=2pt, topsep=2pt}
\setlist[enumerate]{leftmargin=*, itemsep=2pt, topsep=2pt}

\usepackage[most]{tcolorbox}
\usepackage{booktabs}
\usepackage{tikz}
\usetikzlibrary{decorations.pathmorphing, positioning, shapes.geometric, arrows.meta}

% Merkkästchen — wichtige Definition / Kernpunkt
\newtcolorbox{merksatz}[1][]{%
  enhanced, breakable,
  colback=bfwBgTeal,
  colframe=bfwPrimary,
  boxrule=0.6pt,
  arc=4pt,
  left=10pt, right=10pt, top=8pt, bottom=8pt,
  fonttitle=\bfseries\color{bfwPrimaryDark},
  title={\faLightbulb\ Merksatz},
  #1
}

% Eselsbrücke — Mnemoniks
\newtcolorbox{eselsbruecke}[1][]{%
  enhanced, breakable,
  colback=bfwBgRose,
  colframe=bfwAccent,
  boxrule=0.6pt,
  arc=4pt,
  left=10pt, right=10pt, top=8pt, bottom=8pt,
  fonttitle=\bfseries\color{bfwWarn},
  title={Eselsbrücke},
  #1
}

% Beispiel-Box
\newtcolorbox{beispiel}[1][]{%
  enhanced, breakable,
  colback=bfwBgSoft,
  colframe=bfwInkSoft,
  boxrule=0.4pt,
  arc=3pt,
  left=10pt, right=10pt, top=8pt, bottom=8pt,
  fonttitle=\bfseries\color{bfwInk},
  title={Beispiel},
  #1
}

% Lösungs-Box (für Aufgabenhefte)
\newtcolorbox{loesung}[1][]{%
  enhanced, breakable,
  colback=bfwBgAmber!40,
  colframe=bfwSuccess,
  boxrule=0.6pt,
  arc=3pt,
  left=10pt, right=10pt, top=8pt, bottom=8pt,
  fonttitle=\bfseries\color{bfwSuccess!70!black},
  title={Lösung},
  #1
}

% Glossar-Eintrag
\newcommand{\glossareintrag}[2]{%
  \par\noindent\textbf{\color{bfwPrimaryDark}#1} \enspace #2\par\smallskip
}

% Section-Stil — etwas farbiger als Standard
\usepackage{titlesec}
\titleformat{\section}
  {\Large\bfseries\color{bfwPrimaryDark}}
  {\thesection}{0.6em}{}
\titleformat{\subsection}
  {\large\bfseries\color{bfwPrimary}}
  {\thesubsection}{0.5em}{}
\titleformat{\subsubsection}
  {\normalsize\bfseries\color{bfwInk}}
  {\thesubsubsection}{0.4em}{}

% TikZ-Stil "skizzenhaft" — etwas weicher als Rough.js, aber stabil
\tikzset{
  roughbox/.style = {
    draw=bfwPrimaryDark, thick, fill=bfwBgTeal,
    rounded corners=4pt, inner sep=6pt
  },
  roughaccent/.style = {
    draw=bfwAccent, thick, fill=bfwBgAmber,
    rounded corners=4pt, inner sep=6pt
  },
  roughline/.style = {
    draw=bfwInkSoft, thick, -{Stealth[length=2.5mm]}
  }
}

% Bullet-Symbole (bewusst kein FontAwesome — vermeidet Font-Installations-Probleme)
\newcommand{\faLightbulb}{$\bullet$}
\newcommand{\faBookmark}{$\bullet$}

% Header/Footer
\usepackage{fancyhdr}
\pagestyle{fancy}
\fancyhf{}
\renewcommand{\headrulewidth}{0pt}
\fancyfoot[L]{\small\color{bfwInkSoft}\thepage}
\fancyfoot[R]{\small\color{bfwInkSoft} BFW M-064 Beschaffung im Büromanagement}


\title{\vspace*{-1.5cm}\Huge\bfseries\color{bfwPrimaryDark}Aufgabenheft Tag~5\\[0.4em]
  \large\color{bfwInkSoft}\textnormal{Zeitmanagement \& Aufgabenplanung --- Übungen im IHK-Klausurformat}}
\author{\textsf{Büroprozesse gestalten und Arbeitsvorgänge organisieren \textperiodcentered{} M-067}\\
\textsf{\small\copyright\ Can Siebert\,\textperiodcentered\,2026}}
\date{Selbstlernmaterial \textperiodcentered{} 03.07.2026}

\begin{document}
\maketitle
\thispagestyle{empty}

\begin{merksatz}[title={\bfseries So nutzen Sie dieses Aufgabenheft}]
Bearbeiten Sie alle Aufgaben zunächst \emph{ohne} in die Lösungen zu schauen. Schreiben
Sie Ihre Antworten in vollständigen Sätzen --- die IHK-Prüfung verlangt formulierte
Begründungen, nicht bloße Stichworte. Beachten Sie die \emph{Operatoren} (nennen,
erläutern, beurteilen, begründen, anwenden) --- sie geben den geforderten Antwort-Tiefegrad
vor. Erst danach öffnen Sie den Lösungsteil ab Seite~\pageref{loesungen} und vergleichen.
\end{merksatz}

\tableofcontents
\vspace{1em}
\hrule

\section{Aufgaben}

\subsection*{Aufgabe~1 --- Zeitmanagement und Zeitdiebe (10 Punkte)}

\noindent\textbf{\color{bfwAccent}YouTube-Vorbereitung:}\enspace \href{https://www.youtube.com/watch?v=ybspTTfyw8M}{Zeitmanagement --- Grundlagen} (\url{https://www.youtube.com/watch?v=ybspTTfyw8M})

\medskip
Sie arbeiten in der Verwaltung der \emph{Duisdorfer Büro-Konzept~KG}. Ihre neue Kollegin
klagt, sie komme „nie zu ihren eigentlichen Aufgaben”.

\begin{enumerate}[label=(\alph*)]
  \item \textbf{Erläutern} Sie, worum es beim Zeitmanagement grundsätzlich geht. \hfill (2~P.)
  \item \textbf{Nennen} Sie zwei Ziele, die mit dem Zeitmanagement angestrebt werden. \hfill (2~P.)
  \item \textbf{Unterscheiden} Sie äußere und innere Störungen und nennen Sie je zwei Beispiele. \hfill (4~P.)
  \item \textbf{Beurteilen} Sie die Aussage: \emph{„Planung kostet nur Zeit, die ich besser zum Arbeiten nutze.”} \hfill (2~P.)
\end{enumerate}

\subsection*{Aufgabe~2 --- Sägeblatteffekt und stille Stunde (10 Punkte)}

\noindent\textbf{\color{bfwAccent}YouTube-Vorbereitung:}\enspace \href{https://www.youtube.com/watch?v=IUeuEXAWJVg}{Der Sägeblatteffekt} (\url{https://www.youtube.com/watch?v=IUeuEXAWJVg})

\medskip
\begin{enumerate}[label=(\alph*)]
  \item \textbf{Beschreiben} Sie, wie der Sägeblatteffekt entsteht. \hfill (3~P.)
  \item \textbf{Nennen} Sie, wie viel Prozent der Zeit durch Störungen verloren gehen können. \hfill (1~P.)
  \item \textbf{Erläutern} Sie das Konzept der „stillen Stunde”. \hfill (3~P.)
  \item \textbf{Beurteilen} Sie, warum man die stille Stunde nicht zur E-Mail-Kontrolle nutzen sollte. \hfill (3~P.)
\end{enumerate}

\subsection*{Aufgabe~3 --- Leistungskurve und Pausen (12 Punkte)}

\noindent\textbf{\color{bfwAccent}YouTube-Vorbereitung:}\enspace \href{https://www.youtube.com/watch?v=BzziXcCI8p0}{Leistungskurve \& Biorhythmus} (\url{https://www.youtube.com/watch?v=BzziXcCI8p0})

\medskip
\begin{enumerate}[label=(\alph*)]
  \item \textbf{Benennen} Sie das Leistungshoch und das Leistungstief eines Durchschnittsmenschen. \hfill (3~P.)
  \item \textbf{Erläutern} Sie, welche Aufgaben an den Vormittag und welche an den Nachmittag gehören. \hfill (3~P.)
  \item \textbf{Unterscheiden} Sie Morgen- und Abendmensch. \hfill (2~P.)
  \item \textbf{Erklären} Sie die Pausenregelung des §~4 ArbZG bei mehr als sechs bis neun bzw.\ mehr als neun Stunden Arbeitszeit. \hfill (4~P.)
\end{enumerate}

\subsection*{Aufgabe~4 --- ALPEN-Tagesplan erstellen (16 Punkte)}

\noindent\textbf{\color{bfwAccent}YouTube-Vorbereitung:}\enspace \href{https://www.youtube.com/watch?v=j3Q-IedSVtE}{Die ALPEN-Methode} (\url{https://www.youtube.com/watch?v=j3Q-IedSVtE})

\medskip
Frau Albers aus der Ausstattungsberatung hat für heute u.\,a.\ folgende Aufgaben: Termin
beim Werksarzt vereinbaren, Angebot für die Frohn~KG erstellen, neue Prospekte bestellen,
Blumen gießen, Anfrage der Neukunden bearbeiten, Personalabteilung wegen
Weiterbildungszertifikat anrufen.

\begin{enumerate}[label=(\alph*)]
  \item \textbf{Nennen} Sie die fünf Schritte der ALPEN-Methode und erklären Sie sie jeweils kurz. \hfill (5~P.)
  \item \textbf{Wenden} Sie die ALPEN-Methode auf Frau Albers' Aufgaben an: Erstellen Sie einen kurzen Tagesplan mit geschätzter Dauer und Reihenfolge. \hfill (6~P.)
  \item \textbf{Begründen} Sie, warum man nur ca.\ 60~\% der Zeit verplanen sollte. \hfill (3~P.)
  \item \textbf{Beurteilen} Sie, welche der Aufgaben Frau Albers delegieren könnte. \hfill (2~P.)
\end{enumerate}

\subsection*{Aufgabe~5 --- Eisenhower-Prinzip anwenden (14 Punkte)}

\noindent\textbf{\color{bfwAccent}YouTube-Vorbereitung:}\enspace \href{https://www.youtube.com/watch?v=X0Vk4cTyYtE}{Das Eisenhower-Prinzip} (\url{https://www.youtube.com/watch?v=X0Vk4cTyYtE})

\medskip
Ordnen Sie folgende Aufgaben in die vier Felder der Eisenhower-Matrix ein und nennen Sie
die jeweils richtige Handlung:
\begin{enumerate}[label=\arabic*., leftmargin=*]
  \item Die Präsentation für den morgigen Kundentermin ist noch nicht fertig.
  \item Ein Mitarbeiter bittet um Hilfe bei einer Software-Neuinstallation.
  \item Die langfristige Jahresplanung der Abteilung soll vorbereitet werden.
  \item Ein Werbeprospekt eines unbekannten Anbieters liegt auf dem Schreibtisch.
\end{enumerate}

\begin{enumerate}[label=(\alph*)]
  \item \textbf{Erläutern} Sie die beiden Kriterien des Eisenhower-Prinzips und welches Vorrang hat. \hfill (3~P.)
  \item \textbf{Ordnen} Sie die vier Aufgaben den Feldern 1--4 \textbf{zu} und nennen Sie die Handlung. \hfill (8~P.)
  \item \textbf{Beurteilen} Sie, warum „Wichtigkeit vor Dringlichkeit” gehen sollte. \hfill (3~P.)
\end{enumerate}

\subsection*{Aufgabe~6 --- Pareto, ABC, To-do und SMART (14 Punkte)}

\noindent\textbf{\color{bfwAccent}YouTube-Vorbereitung:}\enspace \href{https://www.youtube.com/watch?v=yaVkudeTUKo}{Pareto, ABC \& SMART} (\url{https://www.youtube.com/watch?v=yaVkudeTUKo})

\medskip
\begin{enumerate}[label=(\alph*)]
  \item \textbf{Erläutern} Sie das Pareto-Prinzip (80/20) am Beispiel des Zeitmanagements. \hfill (3~P.)
  \item \textbf{Ordnen} Sie den Kategorien A, B, C der ABC-Analyse Priorität und Handlung \textbf{zu}. \hfill (3~P.)
  \item \textbf{Erläutern} Sie den Unterschied zwischen Checkliste und To-do-Liste. \hfill (3~P.)
  \item \textbf{Wenden} Sie die SMART-Regel an: Formulieren Sie ein vollständiges SMART-Ziel und benennen Sie die fünf Kriterien. \hfill (5~P.)
\end{enumerate}

\vspace{1em}
\noindent\textbf{Punkte gesamt: 76}

\newpage

\section{Lösungen}\label{loesungen}

\subsection*{Lösung Aufgabe~1}
\begin{loesung}
\textbf{(a)} Beim Zeitmanagement geht es darum, die persönlich vorhandene Zeit möglichst
effektiv zu nutzen. Dafür stehen Hilfsmittel und Methoden zur Verfügung; gute Planung
gewinnt im Anschluss Zeit für Berufliches und Privates zurück.

\textbf{(b)} Zwei Ziele (Beispiele): Vermeidung von Stress und Überforderung, Steigerung
der persönlichen Zufriedenheit, ausgewogene Work-Life-Balance, genügend Zeit für die
wesentlichen Aufgaben.

\textbf{(c)} \emph{Innere} Störungen verursacht die Person selbst, z.\,B.\ Lustlosigkeit,
fehlende Selbstdisziplin, Unfähigkeit „nein” zu sagen, Müdigkeit, Tagträumen.
\emph{Äußere} Störungen kommen von anderen Personen oder den Gegebenheiten, z.\,B.\
unordentlicher Arbeitsplatz, Unterbrechung durch Kollegen, klingelnde Telefone, Lärm,
technische Störungen, Wartezeiten am Gemeinschaftsdrucker.

\textbf{(d)} Die Aussage ist falsch. Für die Planung muss man zwar etwas Zeit
investieren, doch bei guter Planung gewinnt man danach deutlich mehr Zeit zurück, weil
Zeitdiebe vermieden, Prioritäten gesetzt und Doppelarbeit verhindert werden.
\end{loesung}

\subsection*{Lösung Aufgabe~2}
\begin{loesung}
\textbf{(a)} Mit jeder Unterbrechung/Störung fällt die Konzentration und
Leistungsfähigkeit ab. Danach braucht man wieder Zeit, um die Konzentration aufzubauen.
Bei häufigen Störungen ergibt sich grafisch ein sägeblattförmiger Verlauf --- der
Sägeblatteffekt.

\textbf{(b)} Bis zu \textbf{28~\%} der Zeit können verloren gehen.

\textbf{(c)} Die stille Stunde ist ein fest eingeplanter Termin mit sich selbst für
ungestörtes, konzentriertes Arbeiten an wichtigen Aufgaben. Man trägt sie in den
Terminplan ein, stellt das Telefon um oder schaltet den Anrufbeantworter ein und legt
sie in eine ohnehin störungsarme Zeit.

\textbf{(d)} Die stille Stunde ist kostbare, störungsfreie Zeit und sollte daher für
wirklich wichtige (A-)Aufgaben genutzt werden. E-Mail-Kontrolle ist meist eine
unwichtige Routine-/C-Aufgabe; sie dafür zu „verbrauchen” würde den eigentlichen Zweck
verfehlen.
\end{loesung}

\subsection*{Lösung Aufgabe~3}
\begin{loesung}
\textbf{(a)} Das Leistungshoch liegt zwischen \textbf{08:00 und 12:00 Uhr}; das größte
Tief kommt \textbf{nach der Mittagspause}.

\textbf{(b)} An den \emph{Vormittag} gehören alle Aufgaben, für die volle
Leistungsfähigkeit nötig ist (wichtige A-Aufgaben); an den \emph{Nachmittag} eignen sich
Routineaufgaben wie die Ablage von Unterlagen.

\textbf{(c)} Der \emph{Morgenmensch} hat seine höchste Leistungsfähigkeit am frühen
Vormittag, der \emph{Abendmensch} am Nachmittag. Beide sollten wichtige Aufgaben in ihre
persönliche Hochphase legen.

\textbf{(d)} §~4 ArbZG: mindestens \textbf{30~Minuten} Ruhepause bei einer Arbeitszeit
von mehr als sechs bis zu neun Stunden, mindestens \textbf{45~Minuten} bei mehr als neun
Stunden. Die Pausen können in Abschnitte von je mindestens 15~Minuten aufgeteilt werden;
länger als sechs Stunden hintereinander darf nicht ohne Pause gearbeitet werden.
\end{loesung}

\subsection*{Lösung Aufgabe~4}
\begin{loesung}
\textbf{(a)} \textbf{A}ufgaben/Termine aufschreiben; \textbf{L}änge der Tätigkeiten
realistisch abschätzen; \textbf{P}ufferzeiten einplanen; \textbf{E}ntscheidungen über
Prioritäten treffen, evtl.\ delegieren; \textbf{N}achkontrolle und nicht Erledigtes neu
einplanen.

\textbf{(b)} Beispielhafter Tagesplan (Reihenfolge nach Priorität und Leistungskurve):
\begin{itemize}[itemsep=0pt]
  \item Vormittag (Hochphase): Angebot für die Frohn~KG erstellen (ca.\ 60~Min, A); Anfrage der Neukunden bearbeiten (ca.\ 45~Min, A).
  \item Mittag/organisatorisch: Termin beim Werksarzt vereinbaren (ca.\ 10~Min); Personalabteilung anrufen (ca.\ 10~Min).
  \item Nachmittag (Routine): neue Prospekte bestellen (ca.\ 15~Min, C); Blumen gießen (ca.\ 5~Min, C).
\end{itemize}
Wichtige Aufgaben liegen am Vormittag, Routine am Nachmittag; zwischen den Blöcken
Pufferzeiten.

\textbf{(c)} Nur ca.\ 60~\% verplanen, damit ca.\ 40~\% als Reserve für Unvorhergesehenes
bleiben. Sonst gerät der Plan schon bei kleinen zeitlichen Abweichungen durcheinander.

\textbf{(d)} Delegierbar sind unwichtige Routineaufgaben, z.\,B.\ das Bestellen der
Prospekte oder das Blumengießen; die wichtigen Aufgaben (Angebot, Neukundenanfrage)
sollte Frau Albers selbst erledigen.
\end{loesung}

\subsection*{Lösung Aufgabe~5}
\begin{loesung}
\textbf{(a)} Die zwei Kriterien sind \emph{Dringlichkeit} (konkreter Termin) und
\emph{Wichtigkeit} (bringt einem Ziel näher, Arbeitserfolg, oft langfristig).
\textbf{Wichtigkeit geht vor Dringlichkeit.}

\textbf{(b)} Zuordnung:
\begin{enumerate}[label=\arabic*., itemsep=0pt]
  \item Präsentation für morgen: wichtig + dringend $\rightarrow$ \textbf{Feld~1}: sofort und selbst erledigen.
  \item Software-Neuinstallation: dringend + nicht wichtig $\rightarrow$ \textbf{Feld~3}: delegieren.
  \item Langfristige Jahresplanung: wichtig + nicht dringend $\rightarrow$ \textbf{Feld~2}: Termin vereinbaren/Terminplanung.
  \item Werbeprospekt unbekannter Anbieter: nicht wichtig + nicht dringend $\rightarrow$ \textbf{Feld~4}: Papierkorb.
\end{enumerate}

\textbf{(c)} Wichtige Aufgaben tragen am stärksten zur Zielerreichung bei, sind aber oft
langfristig und ohne akuten Termin --- sie drohen daher im hektischen Tagesgeschäft
unterzugehen. Wer immer nur dem Dringenden hinterherläuft, vernachlässigt das wirklich
Wichtige; deshalb hat die Wichtigkeit Vorrang.
\end{loesung}

\subsection*{Lösung Aufgabe~6}
\begin{loesung}
\textbf{(a)} Pareto-Prinzip (80/20): 20~\% der eingesetzten Zeit bewirken 80~\% des
Erfolgs; die restlichen 80~\% der Zeit ergeben nur 20~\% des Erfolgs. Konsequenz:
Prioritäten setzen, Wichtiges zuerst, nicht alles mit gleicher Aufmerksamkeit zu 100~\%
erledigen wollen.

\textbf{(b)} A-Aufgaben = hohe Priorität (Muss) $\rightarrow$ morgens, selbst erledigen,
nicht delegieren; B-Aufgaben = mittlere Priorität (Soll) $\rightarrow$ nach den
A-Aufgaben, meist nachmittags; C-Aufgaben = niedrige Priorität (Kann) $\rightarrow$
Routine, geringster Erfolgsbeitrag, ggf.\ delegieren.

\textbf{(c)} Die \emph{To-do-Liste} hält zukünftige Aufgaben fest und organisiert sie
(„Was ist zu tun? Wer macht was bis wann?”) --- zukunftsbezogen. Die \emph{Checkliste}
prüft, ob alle Arbeitsschritte erledigt wurden, und sichert Vollständigkeit/Reihenfolge
--- vergangenheitsbezogen.

\textbf{(d)} SMART = \textbf{S}pezifisch, \textbf{M}essbar, \textbf{A}kzeptiert,
\textbf{R}ealistisch, \textbf{T}erminiert. Beispiel-Ziel: „Verringerung des
Druckerpapierverbrauchs in der Abteilung Beschaffung um 10~\% bis zum 31.12.” --- es ist
spezifisch (Druckerpapier, Abteilung), messbar (10~\%), akzeptiert (positiv,
umsetzbar), realistisch (erreichbar) und terminiert (31.12.).
\end{loesung}

\end{document}
